\chapter{Resultados e Discussões}\label{cap:resultados}
\glsresetall

% =====================================================================
% CAPITULO AGUARDANDO COLETA (estrutura fixada; dados a preencher).
% Alinhado ao desenho de 4 condicoes (C1..C4) consolidado em 2026-07-30.
% As tabelas sao gabaritos com celulas [--] a completar apos a coleta.
% Ordem de apresentacao: amostra -> qualidade de sinal -> verificacoes ->
% primarios -> secundarios -> vereditos -> discussao.
% =====================================================================

Este capítulo apresentará e discutirá os resultados empíricos do estudo, estabelecendo o veredito das hipóteses H1 a H6 formuladas na Seção~\ref{sec:hipoteses}. No momento da redação, o estudo encontra-se em fase de estruturação metodológica e de aprovação ética, e as sessões práticas ainda não foram iniciadas, de modo que não há dados de coleta disponíveis. As seções a seguir fixam a estrutura de apresentação, que segue a ordem recomendada para desenhos com desfechos declarados: caracterização da amostra, qualidade dos sinais, verificação das manipulações, desfechos primários, desfechos secundários e, por fim, o veredito das hipóteses.

% =====================================================================
\section{Caracterização da Amostra}\label{sec:res-amostra}

\ph{descrever a amostra coletada, com o número de completadores por sequência do quadrado de Williams, idade, sexo ou gênero, experiência prévia com \gls{rv}, jogos de ritmo, tênis de mesa, meditação e música \textit{lo-fi}, e escores basais de \gls{pss} e \gls{gad}. Relatar o fluxo de participantes desde o recrutamento, com perdas e exclusões discriminadas por motivo: desistência, falha de sensor, artefato excessivo e interrupção por cybersickness.}

% =====================================================================
\section{Qualidade dos Dados Fisiológicos}\label{sec:res-qualidade}

\ph{reportar, conforme exigido pelas recomendações de publicação adotadas, o modelo e a versão de firmware dos dispositivos, a frequência de amostragem, o tipo de eletrodo e o meio de contato, a temperatura e a umidade da sala, o percentual de intervalos RR corrigidos por participante e condição, e o percentual de janelas rejeitadas por artefato \citep{Laborde2017,Boucsein2012}.}

% =====================================================================
\section{Verificações de Manipulação}\label{sec:res-manipulacao}

Antes de qualquer teste de hipótese, verificam-se as manipulações. Um resultado nulo em qualquer delas condiciona a interpretação de todos os desfechos subsequentes.

\ph{relatar: (i) se a indução de estresse elevou a VAS de ansiedade e estresse e a ativação do SAM, em cada uma das quatro sessões, e se essa elevação se atenuou ao longo delas, indicando habituação; (ii) se a carga do NASA-TLX e o esforço de Borg foram maiores em C2 do que em C1; (iii) se a apreciação da trilha foi equivalente entre C1 e C3; (iv) se o erro objetivo de sincronização permaneceu baixo em C1 e C2; (v) se a cybersickness se manteve em níveis aceitáveis nas condições em RV.}

\begin{table}[H]
    \centering
    \caption{Verificações de manipulação, a preencher após a coleta com médias, desvios, estatística de teste e valor de $p$.}
    \label{tab:res-manipulacao}
    \small
    \begin{tabularx}{\textwidth}{p{5.0cm}YY}
        \toprule
        \textbf{Verificação} & \textbf{Resultado} & \textbf{Conclusão} \\
        \midrule
        Estressor elevou \gls{vas}-Ansiedade & [--] & [--] \\
        Atenuação ao longo das sessões & [--] & [--] \\
        \gls{tlx}: C2 $>$ C1 & [--] & [--] \\
        Borg: C2 $>$ C1 $>$ C3, C4 & [--] & [--] \\
        Trilha igualmente apreciada (C1, C3) & [--] & [--] \\
        Erro de sincronização baixo (C1, C2) & [--] & [--] \\
        \bottomrule
    \end{tabularx}\\[2pt]
    {\footnotesize Fonte: Elaborada pelo autor.}
\end{table}

% =====================================================================
\section{Desfecho Psicológico Primário}\label{sec:res-psicologico}

\ph{apresentar a variação da VAS de ansiedade entre a medida pós-estressor e a pós-intervenção, por condição, e os três contrastes planejados com correção de Holm, com estimativas, intervalos de confiança de 95\% e $d_z$.}

\begin{table}[H]
    \centering
    \caption{Desfecho psicológico primário: $\Delta$ \gls{vas}-Ansiedade (pós-intervenção menos pós-estressor), por condição. Valores mais negativos indicam maior recuperação.}
    \label{tab:res-primario}
    \small
    \begin{tabularx}{\textwidth}{p{3.2cm}YYYY}
        \toprule
        \textbf{Desfecho} & \textbf{C1 Relaxante} & \textbf{C2 Animado} & \textbf{C3 Trilha} & \textbf{C4 Meditação} \\
        \midrule
        $\Delta$ \gls{vas}-Ansiedade & [--] & [--] & [--] & [--] \\
        \bottomrule
    \end{tabularx}\\[2pt]
    {\footnotesize Fonte: Elaborada pelo autor.}
\end{table}

\begin{table}[H]
    \centering
    \caption{Contrastes planejados sobre o desfecho psicológico primário, com correção de Holm.}
    \label{tab:res-contrastes}
    \small
    \begin{tabularx}{\textwidth}{p{2.0cm}p{4.2cm}YY}
        \toprule
        \textbf{Contraste} & \textbf{Hipótese associada} & \textbf{Estimativa (IC $95\%$)} & \textbf{$d_z$, $p$ ajustado} \\
        \midrule
        C1 \textit{vs} C3 & H1 (ganho sobre a música) & [--] & [--] \\
        C1 \textit{vs} C2 & H2 (projeto relaxante) & [--] & [--] \\
        C1 \textit{vs} C4 & H3 (não inferioridade) & [--] & [--] \\
        \bottomrule
    \end{tabularx}\\[2pt]
    {\footnotesize Fonte: Elaborada pelo autor.}
\end{table}

% =====================================================================
\section{Desfechos Fisiológicos}\label{sec:res-fisiologico}

\ph{apresentar a variação do logaritmo da RMSSD entre linha de base e recuperação, desfecho fisiológico primário, e os desfechos secundários de SCL, FC média e SDNN. Reportar os modelos com e sem o esforço percebido como covariável, conforme a análise de sensibilidade prevista na Seção~\ref{sec:analise}.}

\begin{table}[H]
    \centering
    \caption{Desfechos fisiológicos por condição, a preencher com médias, desvios e intervalos de confiança.}
    \label{tab:res-fisio}
    \small
    \begin{tabularx}{\textwidth}{p{4.0cm}YYYY}
        \toprule
        \textbf{Desfecho} & \textbf{C1} & \textbf{C2} & \textbf{C3} & \textbf{C4} \\
        \midrule
        $\Delta \ln(\mathrm{RMSSD})$ (recup. $-$ base) & [--] & [--] & [--] & [--] \\
        $\Delta$ \gls{scl} (recup. $-$ pós-estressor) & [--] & [--] & [--] & [--] \\
        \gls{fc} média na recuperação & [--] & [--] & [--] & [--] \\
        \gls{sdnn} na recuperação & [--] & [--] & [--] & [--] \\
        \bottomrule
    \end{tabularx}\\[2pt]
    {\footnotesize Fonte: Elaborada pelo autor.}
\end{table}

% =====================================================================
\section{Desfechos Secundários e de Experiência}\label{sec:res-secundarios}

\ph{reportar, por condição, a VAS de estresse e calma, o SAM, a subescala de tensão do BRUMS, o NASA-TLX, o Borg CR10 e a fruição da atividade. Reportar separadamente, como caracterização das condições em que se aplicam, o IPQ, o SSQ, a sincronização percebida e o desempenho no jogo. Incluir a análise qualitativa das respostas abertas sobre satisfação, frustração e sugestões.}

% =====================================================================
\section{Veredito das Hipóteses}\label{sec:res-vereditos}

\ph{preencher o veredito de cada hipótese, como sustentada, parcialmente sustentada ou refutada, com a evidência correspondente.}

\begin{table}[H]
    \centering
    \caption{Síntese dos vereditos das hipóteses H1 a H6, a completar após a análise.}
    \label{tab:res-vereditos}
    \small
    \begin{tabularx}{\textwidth}{p{1.0cm}Xp{3.0cm}}
        \toprule
        \textbf{Hip.} & \textbf{O que afirma} & \textbf{Veredito} \\
        \midrule
        H1 & C1 supera C3: o jogo acrescenta recuperação à da música & [--] \\
        H2 & C1 supera C2: o projeto concebido para acalmar importa & [--] \\
        H3 & C1 não é inferior a C4, a meditação guiada em \gls{rv} & [--] \\
        H4 & C2 impõe maior carga e esforço que C1 & [--] \\
        H5 & C1 supera C2 e C3 na recuperação eletrodérmica & [--] \\
        H6 & Convergência entre medidas subjetivas e fisiológicas & [--] \\
        \bottomrule
    \end{tabularx}\\[2pt]
    {\footnotesize Fonte: Elaborada pelo autor.}
\end{table}

% =====================================================================
\section{Discussão}\label{sec:discussao}

\ph{discutir os vereditos, conectar os achados à literatura da Seção~\ref{sec:trabalhos-relacionados} e distinguir efeitos psicológicos, fisiológicos e de experiência.}

Os cenários interpretativos abaixo foram previstos antes da coleta e devem orientar a discussão.

\begin{itemize}
    \item \textbf{C1 supera C3 e C2.} Cenário que sustenta integralmente a proposta: o jogo acrescenta recuperação à da música, e essa vantagem depende de a experiência ter sido projetada para acalmar. É o resultado que fundamenta recomendações de projeto.
    \item \textbf{C1 supera C3, mas equipara-se a C2.} O benefício viria de jogar em \gls{rv} com música, independentemente do caráter da experiência, o que enfraqueceria a tese do projeto relaxante e sugeriria que o engajamento, por si, explica o efeito. Convergiria com \citet{RubioArias2024}, que atribuem a vantagem dos \textit{exergames} à fruição, e não à intensidade.
    \item \textbf{C1 equipara-se a C3.} O benefício seria atribuível sobretudo à trilha musical, resultado compatível com \citet{Dsouza2024}. Dado o caráter de comparador ativo de C3 (Seção~\ref{sec:limitacoes}), isso não indicaria ineficácia, mas ausência de ganho sobre ouvir a mesma música.
    \item \textbf{C1 não inferior a C4.} Resultado de maior valor prático: posicionaria um \textit{exergame} rítmico como alternativa à meditação guiada, com o argumento adicional de maior fruição, caso confirmado pelos desfechos de experiência. Dialoga diretamente com \citet{Ahn2025}, que observaram vantagem da \gls{rv} sobre técnica convencional em desfecho subjetivo.
    \item \textbf{Divergência entre canais.} Melhora subjetiva sem recuperação fisiológica, ou o inverso, padrão já observado por \citet{Kim2021} ao comparar técnicas de relaxamento. Nesse caso, ponderar o esforço físico desigual entre condições, a sensibilidade limitada da \gls{eda} e possíveis efeitos de expectativa (Seção~\ref{sec:limitacoes}).
\end{itemize}

\ph{encerrar a discussão relacionando os achados às recomendações de projeto de jogos de bem-estar que deles decorrem.}
